\documentclass[12pt]{article}
\usepackage{geometry}
\geometry{letterpaper, margin=1in}
\usepackage{hyperref}
\usepackage[usenames,dvipsnames]{xcolor}
\usepackage{microtype}

\definecolor{accent}{RGB}{26, 60, 110}   % matches resume navy

\begin{document}

%----------HEADER----------
\begin{center}
  {\LARGE\textbf{\textcolor{accent}{SITONG LI}}} \\[4pt]
  \small Minneapolis, MN
  \;\textbullet\;
  \href{mailto:ruc.sitongli@gmail.com}{ruc.sitongli@gmail.com}
  \;\textbullet\;
  (952) 201-4026
\end{center}

\vspace{1em}

\noindent \today

\vspace{1em}

\noindent Dear Professor Ashokkumar,

\vspace{1em}

I am writing to apply for the Research Assistant position at the We-Search Lab at Harvard University. Your lab's focus on the psychology of political discourse---how language reflects and shapes identity processes and group dynamics---sits precisely at the intersection of questions I have been pursuing through my own computational research. I am eager to bring my experience building large-scale natural language processing pipelines to projects that illuminate the social psychology of how people communicate and position themselves in political life.

My deepest relevant experience comes from a research assistantship with Professor William Cassidy at Washington University in St.\ Louis, where I analyzed natural language datasets drawn from 1950s congressional hearings to study partisan identity and ideological divergence. I developed and maintained the full text analysis pipeline: digitizing scanned PDFs using a multimodal LLM, classifying topics and extracting policy indicators with GPT-4, and resolving speaker identities and linking them to party affiliation via BERT-based entity resolution. To measure how group membership shapes rhetoric, I built a word-embedding model inspired by Gentzkow, Shapiro, and Taddy (2019) to quantify the degree to which language diverges across party lines across sessions. The core questions in that project---how political identity is expressed through language, how in-group versus out-group dynamics manifest in discourse, and how these patterns change over time---map directly onto the identity and group dynamics questions the We-Search Lab is designed to answer. I would welcome the chance to bring those same methods and instincts to social media data and other natural language corpora at your lab.

Beyond political text, I have built scalable data collection and analysis pipelines across several projects. For Professor Haiwen Zhang at the University of Minnesota, I collected and processed over 100,000 regulatory emails and 10,000 SEC filings, building OCR preprocessing workflows, applying GPT-4 for large-scale topic and entity extraction, and validating findings using regression models in Stata and Python. These projects gave me practical experience with the full research data lifecycle: from raw data ingestion and cleaning, through NLP-based feature construction, to statistical analysis and manuscript-ready outputs. I am proficient in Python and R and have hands-on experience with a wide range of NLP methods, including transformer-based classifiers, large language model prompting, and word embedding models. I take reproducibility seriously and write modular, documented code that collaborators can build on.

I am highly motivated to work in an academic research environment and am committed to contributing to every stage of the research process---from dataset construction and pipeline maintenance to collaborative project development and manuscript writing. I am available to begin immediately and am committed to the full two-year term.

My two references, who can speak to my research skills and work ethic, are:

\begin{itemize}
  \item \textbf{William M.\ Cassidy}, Assistant Professor of Finance, Washington University in St.\ Louis --- \href{mailto:cassidyw@wustl.edu}{cassidyw@wustl.edu}, +1\,(225)\,368-7616
  \item \textbf{Haiwen (Helen) Zhang}, Professor of Accounting, University of Minnesota --- \href{mailto:zhan0400@umn.edu}{zhan0400@umn.edu}, +1\,(612)\,624-9818
\end{itemize}

Thank you for your time and consideration. I would be glad to discuss how my background fits the needs of the We-Search Lab.

\vspace{1em}

\noindent Sincerely,\\[0.5em]
\textbf{Sitong Li}

\end{document}
